\documentclass[11pt,a4paper]{article}

% --- PAKETE & SETUP ---
\usepackage[ngerman]{babel}
\usepackage[margin=2.5cm, top=3cm, bottom=3cm]{geometry}
\usepackage{amsmath, amssymb, amsthm}
\usepackage{graphicx}
\usepackage{booktabs}
\usepackage{tabularx}
\usepackage{array}
\usepackage{xcolor}
\usepackage{fancyhdr}
\usepackage{titlesec}
\usepackage{enumitem}
\usepackage{caption}
\usepackage{framed}
\usepackage{hyperref}
\usepackage{lipsum} % Für Platzhalter-Text

% --- FARBDEFINITIONEN ---
\definecolor{primary}{RGB}{20, 50, 90}       % Dunkelblau
\definecolor{secondary}{RGB}{70, 110, 150}   % Muted Blau

% --- HYPERREF SETUP ---
\hypersetup{
    colorlinks=true,
    linkcolor=primary,
    citecolor=primary,
    urlcolor=secondary
}

% --- MANUELLE PLATZHALTER (HIER IHRE DATEN EINTRAGEN) ---
\newcommand{\thethema}{Santander Customer Transaction Prediction}
\newcommand{\theautor}{Maria Musterfrau}
\newcommand{\thematrikelnr}{1234567}
\newcommand{\thestudiengang}{Data Science \& Analytics}
\newcommand{\themodul}{Advanced Data Analytics}
\newcommand{\thedozent}{Dr. Jan-Matthis Waack}
\newcommand{\thesemester}{SoSe 2026}
\newcommand{\theabgabedatum}{1. September 2026}
\newcommand{\thegruppennummer}{138}
\newcommand{\thesprecherin}{Maria Musterfrau (Sprecherin)}
\newcommand{\themitglieder}{Maria Musterfrau, Max Mustermann, Erika Mustermann}

% Schalter: Ist dieses Portfolio von der Sprecherin?
\newif\ifsprecherin
\sprecherintrue % Auf \sprecherinfalse setzen, falls Sie NICHT Sprecherin sind!

% --- KOPF- UND FUSSZEILEN ---
\pagestyle{fancy}
\fancyhf{}
\fancyhead[L]{\small\color{gray} e-Portfolio | ADA}
\fancyhead[R]{\small\color{gray} Gruppe \thegruppennummer}
\fancyfoot[C]{\thepage}
\renewcommand{\headrulewidth}{0.4pt}
\renewcommand{\footrulewidth}{0.4pt}

% --- SEKTIONSFORMATIERUNG ---
\titleformat{\section}
  {\normalfont\Large\bfseries\color{primary}}{\thesection}{1em}{}[\titlerule]
\titleformat{\subsection}
  {\normalfont\large\bfseries\color{secondary}}{\thesubsection}{1em}{}

% ==============================================================================
\begin{document}

% ------------------------------------------------------------------------------
% POS 1: TITELSEITE (DECKBLATT) - Max. 1 Seite
% ------------------------------------------------------------------------------
\begin{titlepage}
    \thispagestyle{empty}
    \centering
    \vspace*{1cm}
    
    {\Large\scshape\color{secondary} e-Portfolio Ausarbeitung}\\[0.5cm]
    {\huge\bfseries\color{primary} \thethema}\\[1.5cm]
    
    \rule{\linewidth}{1pt}\\[1cm]
    
    \begin{minipage}{0.48\textwidth}
        \begin{flushleft} \large
            \textbf{Autor/in:}\\
            \theautor\\
            Matrikelnr.: \thematrikelnr\\
            \thestudiengang\\[0.5cm]
            
            \textbf{Gruppe:}\\
            Gruppe \thegruppennummer\\[0.2cm]
            \textbf{Sprecherin:}\\
            \thesprecherin
        \end{flushleft}
    \end{minipage}
    \begin{minipage}{0.48\textwidth}
        \begin{flushright} \large
            \textbf{Veranstaltung:}\\
            \themodul\\[0.3cm]
            \textbf{Dozent:}\\
            \thedozent\\[0.3cm]
            \textbf{Semester:}\\
            \thesemester\\[0.3cm]
            \textbf{Abgabedatum:}\\
            \theabgabedatum
        \end{flushright}
    \end{minipage}
    
    \vfill
    
{\large
        \textbf{Gruppenmitglieder:}\\[0.2cm]
        \themitglieder
    }
    
    \vfill
    {\small Modulverantwortlicher: Dr. Jan-Matthis Waack}
\end{titlepage}

\newpage
\pagenumbering{arabic}

% ------------------------------------------------------------------------------
% INHALTSVERZEICHNIS
% ------------------------------------------------------------------------------
\tableofcontents
\newpage

% ------------------------------------------------------------------------------
% POS 3: ZUSAMMENFASSUNG DES LOGBUCHS - Max. 4 Seiten
% ------------------------------------------------------------------------------
\section{Zusammenfassung des Logbuchs (Seite 1 von 4)}
\label{sec:zusammenfassung}

\textit{Hinweis: Ausformulierte Darstellung des Arbeits- und Methodeneinsatzes auf Basis des individuellen Logbuchs.}

\subsection{Teil A: Grundlagen und Arbeitsprozess}
\lipsum[1-3]

\newpage
\subsection*{Zusammenfassung des Logbuchs (Seite 2 von 4)}
\lipsum[4-6]

\newpage
\subsection{Teil B: Reflexion und Einordnung (Seite 3 von 4)}
\lipsum[7-8]

\begin{table}[htbp]
    \centering
    \caption{Ergebnisse des Modellvergleichs}
    \label{tab:ergebnisse}
    \begin{tabular}{lcccc}
        \toprule
        \textbf{Modell} & \textbf{Accuracy} & \textbf{Precision} & \textbf{Recall} & \textbf{F1-Score} \\
        \midrule
        Baseline & 0.75 & 0.70 & 0.68 & 0.69 \\
        Random Forest & 0.86 & 0.83 & 0.81 & 0.82 \\
        XGBoost & 0.90 & 0.88 & 0.87 & 0.87 \\
        \bottomrule
    \end{tabular}
\end{table}

\newpage
\subsection{Teil C: Abbildungen und Tabellen (Seite 4 von 4)}
\lipsum[9-11]

\newpage

% ------------------------------------------------------------------------------
% POS 4: EINSCHÄTZUNG DER GRUPPENMITGLIEDER - Max. 1 Seite
% ------------------------------------------------------------------------------
\section{Einschätzung der Gruppenmitglieder (Max. 1 Seite)}
\label{sec:einschaetzung}

\textit{Hinweis: Vertrauliche, sachliche Einschätzung der Zusammenarbeit. Pro Gruppenmitglied mindestens 1 Stärke und 1 Schwäche benennen.}

\subsection*{Mitglied 1: Max Mustermann}
\begin{itemize}
    \item \textbf{Stärke:} Äußerst strukturiert bei der Erstellung der Daten-Pipelines.
    \item \textbf{Schwäche:} Kommunizierte Zwischenstände in den Meetings manchmal sehr kurzfristig.
\end{itemize}

\subsection*{Mitglied 2: Erika Mustermann}
\begin{itemize}
    \item \textbf{Stärke:} Sehr gründlich bei der Visualisierung der explorativen Datenanalyse.
    \item \textbf{Schwäche:} Neigte dazu, sich in Einzeldetails zu vertiefen.
\end{itemize}

\vfill
\textit{[Platzhalter für weitere Anmerkungen zur Teamarbeit]}

\newpage

% ------------------------------------------------------------------------------
% POS 5: ERGEBNISSE DER GRUPPE - Max. 3 Seiten (Nur Sprecherin)
% ------------------------------------------------------------------------------
\ifsprecherin
\section{Ergebnisse der Gruppe (Seite 1 von 3)}
\label{sec:gruppenergebnisse}

\textit{(Dieser Abschnitt ist nur im e-Portfolio der Sprecherin enthalten.)}

\subsection{Synthese der Ergebnisse}
\lipsum[12-14]

\newpage
\subsection*{Ergebnisse der Gruppe (Seite 2 von 3)}
\lipsum[15-17]

\newpage
\subsection{Diskussion \& Gesamterkenntnis (Seite 3 von 3)}
\lipsum[18-20]

\newpage
\fi

% ------------------------------------------------------------------------------
% POS 6: PROTOKOLL DER GRUPPENTREFFEN - Ca. 1-2 Seiten (Nur Sprecherin)
% ------------------------------------------------------------------------------
\ifsprecherin
\section{Protokoll der Gruppentreffen (Seite 1 von 2)}
\label{sec:protokoll}

\textit{(Dieser Abschnitt ist nur im e-Portfolio der Sprecherin enthalten.)}

\subsection*{Gruppentreffen 1}
\begin{itemize}[leftmargin=*]
    \item \textbf{Datum:} 10. Mai 2026 | \textbf{Teilnehmende:} Alle
    \item \textbf{Besprochene Inhalte:} Wahl der Sprecherin, Verteilung der Vorbereitungsschritte.
    \item \textbf{Vereinbarte Schritte:} Vorbereitung der EDA und Aufsetzen der Ordnerstruktur.
\end{itemize}

\subsection*{Gruppentreffen 2}
\begin{itemize}[leftmargin=*]
    \item \textbf{Datum:} 15. Juni 2026 | \textbf{Teilnehmende:} Alle
    \item \textbf{Besprochene Inhalte:} Evaluation der ersten Modellvarianten.
\end{itemize}

\newpage
\subsection*{Protokoll der Gruppentreffen (Seite 2 von 2)}
\subsection*{Gruppentreffen 3}
\begin{itemize}[leftmargin=*]
    \item \textbf{Datum:} 20. Juli 2026 | \textbf{Teilnehmende:} Alle
    \item \textbf{Besprochene Inhalte:} Finale Vorbereitung des e-Portfolios und Zusammenführung der Ergebnisse.
\end{itemize}
\lipsum[21-22]

\newpage
\fi

% ------------------------------------------------------------------------------
% POS 7: QUELLEN UND HILFSMITTEL - Ca. 1 Seite
% ------------------------------------------------------------------------------
\section{Quellen und Hilfsmittel (Ca. 1 Seite)}
\label{sec:quellen}

\textit{Hinweis: Alle verwendeten Quellen, Bibliotheken und KI-Tools angeben.}

\subsection{Literatur- und Webquellen}
\begin{enumerate}[label={[\arabic*]}]
    \item Pedregosa, F. et al. (2011). Scikit-learn: Machine Learning in Python. \textit{JMLR}, 12, 2825-2830.
    \item XGBoost Documentation (2026). \url{https://xgboost.readthedocs.io/} (Zugriff: 01.08.2026).
\end{enumerate}

\subsection{Einsatz Generativer KI}
\begin{itemize}
    \item \textbf{Verwendetes Tool:} ChatGPT / Claude / Gemini
    \item \textbf{Einsatzbereich:} Unterstützung bei der Erstellung des LaTeX-Layouts und Überprüfung auf Tippfehler.
    \item \textbf{Beispiel-Prompts:} \texttt{"Erstelle eine übersichtliche LaTeX Table für einen Modellvergleich."}
\end{itemize}

\newpage

% ------------------------------------------------------------------------------
% POS 2: LOGBUCH (ANHANG) - Zählt NICHT zum Seitenumfang!
% ------------------------------------------------------------------------------
\appendix
\section{Anhang: Individuelles Logbuch}
\label{sec:logbuch}

\textit{Das Logbuch wird als Anhang beigefügt und zählt nicht zum offiziellen Seitenumfang, wird aber als eigener Bestandteil bewertet. Es belegt den Workload ($\approx 100$ Stunden).}

\subsection*{Logbucheintrag 1}
\label{log:eintrag1}
\begin{itemize}[leftmargin=*]
    \item \textbf{Datum:} 05. Mai 2026 | \textbf{Dauer:} 4 Std.
    \item \textbf{Titel:} Einrichtung der Umgebung \& Erstsichtung
    \item \textbf{Beschreibung:} Aufsetzen des Git-Repositories und Vorbereitung der Datenstruktur.
\end{itemize}

\end{document}
